\chapter{Discussion}
\label{chap:discussion}

The main advantage of embedding-guided $A^*$ lies in its ability to reduce search effort. It avoids the exhaustive exploration of BFS by prioritizing promising nodes and avoids the repeated neural inference of RL by relying on precomputed node embeddings. This makes the approach particularly suitable for large or dense causal graphs, where exhaustive traversal becomes increasingly expensive. This efficiency, however, comes with a trade-off because bounded exploration may terminate before an existing path is discovered, reducing recall. The reliability of the answers also depends on the underlying causal knowledge graph. Automatically constructed resources such as CauseNet and CEG may contain missing, noisy, ambiguous, or overly general relations, so a discovered path represents a connection encoded in the graph rather than definitive evidence of a real-world causal relationship. In addition, the approach assumes that causal relations can be meaningfully composed along a directed path, although causality is not necessarily transitive. Finally, the method is limited to binary causal questions. It determines whether a directed path between the graph nodes corresponding to the cause and effect can be found, but does not assess the quality, plausibility, or strength of the discovered causal relation. It also does not estimate causal effect sizes, compare alternative causes, or support interventional and counterfactual reasoning. These limitations restrict the proposed approach to efficient causal path discovery rather than general causal inference.

